\documentclass{article}

% Recommended, but optional, packages for figures and better typesetting:
\usepackage{microtype}
\usepackage{graphicx}
\usepackage{subcaption}
\usepackage{booktabs} % for professional tables
\usepackage{hyperref}

% Attempt to make hyperref and algorithmic work together better:
\newcommand{\theHalgorithm}{\arabic{algorithm}}
\usepackage{algorithm}
\usepackage{algorithmic}

% Use the following line for the initial blind version submitted for review:
\usepackage{icml2026}

\usepackage{amsmath}
\usepackage{amssymb}
\usepackage{mathtools}
\usepackage{amsthm}

% if you use cleveref..
\usepackage[capitalize,noabbrev]{cleveref}

%%%%%%%%%%%%%%%%%%%%%%%%%%%%%%%%
% THEOREMS
%%%%%%%%%%%%%%%%%%%%%%%%%%%%%%%%
\theoremstyle{plain}
\newtheorem{theorem}{Theorem}[section]
\newtheorem{proposition}[theorem]{Proposition}
\newtheorem{lemma}[theorem]{Lemma}
\newtheorem{corollary}[theorem]{Corollary}
\theoremstyle{definition}
\newtheorem{definition}[theorem]{Definition}
\newtheorem{assumption}[theorem]{Assumption}
\theoremstyle{remark}
\newtheorem{remark}[theorem]{Remark}

\icmltitlerunning{Deconstructing the Localization Illusion: Task-Agnostic Model Merging}

\begin{document}

\twocolumn[
  \icmltitle{Deconstructing the Localization Illusion: Task-Agnostic Multi-Task Model Merging via Self-Routing Activation Calibration}

  \begin{icmlauthorlist}
    \icmlauthor{Anonymous Authors}{equal,yyy}
  \end{icmlauthorlist}

  \icmlaffiliation{yyy}{Department of Computer Science, University of Machine Learning, Location, Country}
  \icmlcorrespondingauthor{Anonymous Authors}{anonymous@uml.edu}

  \icmlkeywords{Model Merging, Activation Calibration, Multi-Task Learning, ResNet}

  \vskip 0.3in
]

\printAffiliationsAndNotice{}  % no special notice

\begin{abstract}
  Multi-task model merging has emerged as a cost-effective and training-free paradigm to consolidate multiple task-specific expert networks into a single multi-task model. However, directly averaging model weights often leads to significant performance degradation due to parameter interference and representation mismatch. Recent works attempt to mitigate this via activation calibration, which scales activations based on task-specific statistics. In this work, we deconstruct activation calibration and challenge the assumption that calibration must be static or require an explicit task identifier at test time. Building on the \emph{Localization Illusion}---which states that representation collapse is highly localized in late layers while early layers remain intact and retain task-specific signatures---we propose \textbf{Self-Routing Activation Calibration (SRAC)}. SRAC leverages early-layer representations to dynamically compute soft, routing similarity weights against task prototypes on the fly, and uses them to dynamically interpolate optimal scale and bias calibration parameters in subsequent layers. Empirically, we demonstrate that SRAC is completely task-agnostic at test time, resolves the \emph{Sparsity Trap}, and recovers near-expert performance, outperforming existing static task-agnostic approaches by a significant margin.
\end{abstract}

\section{Introduction}
\label{sec:intro}

The pretrain-then-finetune paradigm has revolutionized deep learning, enabling the deployment of highly capable models across diverse downstream tasks. While fine-tuning on task-specific datasets yields specialized expert models with exceptional performance, maintaining, storing, and serving dozens of these specialized expert networks simultaneously introduces astronomical computational, storage, and operational overheads. Multi-task model merging has recently emerged as a compelling, training-free alternative: by consolidating the weights of multiple expert models (which share a common pre-trained initialization progenitor) into a single, unified multi-task network, one can perform multiple tasks without additional training or parameter inflation \cite{tiesmerging2023, isonotask2025, della2024, modelingmultitask2025}.

Despite its theoretical and practical appeal, directly averaging the parameter weights of multiple experts (i.e., simple Weight Averaging) typically results in severe performance degradation. This degradation stems from intense parameter interference and representation misalignment in the intermediate layers of the network, which destroys the delicate structures learned during task-specific fine-tuning \cite{tiesmerging2023, whoeverstartedinterference2025}. To alleviate this representation mismatch, activation calibration and representation alignment techniques have been proposed \cite{submission3, submission7}. These methods adapt the intermediate representations of the merged model by adjusting activation statistics (mean and variance) to match those of the individual experts, typically using a tiny calibration dataset of just a few dozen samples.

In this work, we perform a rigorous, systematic deconstruction of activation calibration in multi-task model merging and expose several fundamental limitations and core assumptions of current methodologies. First, we explore the \textbf{Localization Illusion} (first identified in \cite{submission3}), which reveals that representation collapse is not uniform across layers. Centered Kernel Alignment (CKA) analyses show that early and middle layers of merged models remain highly intact (having a CKA similarity $> 0.95$ with the expert models), with the catastrophic representation collapse being heavily localized in the final residual blocks (such as Layer 4 in ResNet-18). This implies that global, uniform activation calibration across all layers is not only redundant but actively harmful, as it injects noise into well-behaved, intact layers.

Second, we uncover the \textbf{Sparsity Trap}: standard channel-wise activation calibration under non-linearities like ReLU is highly unstable. In sparse feature maps, many channels have near-zero or zero activations on small calibration sets. Normalizing by the merged model's standard deviation in these channels leads to division-by-zero or extreme noise amplification, which completely corrupts the representation space and results in random-guessing performance.

Third, we challenge the restrictive assumption that activation calibration must either be static and task-agnostic (like SP-TAAC \cite{submission8}) or require an oracle task-ID at test time (like TCAC \cite{submission3}). To resolve these limitations, we propose \textbf{Self-Routing Activation Calibration (SRAC)}, a novel, training-free, and fully task-agnostic framework that dynamically calibrates intermediate representations on the fly. 

SRAC represents a radical conceptual departure: rather than trying to fix or ignore representation collapse, we turn the Localization Illusion into a powerful routing mechanism. Since early layers are highly intact and preserve distinct task-specific signatures, we extract task prototypes from them. At test time, SRAC computes soft, sample-wise similarity routing weights against these prototypes during the forward pass. These routing weights are then used to dynamically interpolate sparsity-preserving calibration parameters and route final classification heads. Conceptually, SRAC can be viewed as a **zero-shot, training-free, representation-routed Mixture-of-Experts (MoE)**, where the intact early layers act as implicit, highly accurate task routers without needing any backpropagation or gating network training.

Our main contributions are as follows:
\begin{itemize}
  \item We empirically deconstruct the activation calibration landscape, detailing the mechanics of the \emph{Localization Illusion}, the \emph{Sparsity Trap}, and \emph{Head-Backbone Confounding}.
  \item We present \textbf{Self-Routing Activation Calibration (SRAC)}, an elegant dynamic routing approach that is completely task-agnostic at test time, preserves activation sparsity, and dynamically interpolates calibration parameters.
  \item We formulate and evaluate multiple dynamic routing paradigms, including \textbf{Dynamic BatchNorm Parameter Routing (DBPR)} and \textbf{Self-Routing Layer-wise Scaling (SRLS)}, proving their relative strengths and bounds in uncalibrated and calibrated representation spaces.
  \item We conduct exhaustive architectural sweeps, proving that \textbf{Layer 2 is the optimal routing anchor} in the entire network, and show that N-TAAC + SRAC achieves peak task-agnostic performance, recovering 93.6\% of the oracle task-ID performance on a multi-task suite.
\end{itemize}

\section{Related Work}
\label{sec:related}

\textbf{Model Merging in Parameter Space:} 
Model merging aims to unify multiple specialized expert networks fine-tuned from a shared progenitor. Simple Weight Averaging (WA) serves as a baseline, but modern techniques like TIES-Merging \cite{tiesmerging2023} and DARE \cite{whatmatters2024} prune and align parameter updates to reduce task interference. Other works model merging as projective gradient descent \cite{modelingmultitask2025} or identify common and task-specific subspaces to prevent performance degradation \cite{taskleftbehind2025}. Spectral methods, such as Task Singular Vectors \cite{taskvectors2024}, disentangle weight matrices via SVD. These methods operate purely in parameter space and do not adapt to activation-level dynamics at test time.

\textbf{Activation Calibration and Representation Alignment:} 
Rather than modifying weights, representation calibration operates in activation space. Native Joint Calibration (N-TAAC) \cite{submission7} updates BatchNorm statistics on a joint task dataset to adapt the merged network's internal distributions. Task-Conditional Activation Calibration (TCAC) \cite{submission3} performs channel-wise affine scaling of intermediate activations but requires an oracle task-ID at test time and suffers from the Sparsity Trap under ReLU \cite{submission8}. SP-TAAC \cite{submission8} mitigates this via global, layer-wise scaling but remains static and task-agnostic. REDA \cite{submission7} combines representation calibration with decision boundary alignment. Our work bridges the gap by introducing dynamic, input-dependent routing.

\section{Deconstructing Activation Calibration}
\label{sec:deconstruction}

Let $f_{\text{merged}}$ be a merged multi-task model and $f_{\text{exp}}^{(k)}$ be the task-specific expert for task $k$. For a layer $l$ and input $x$, let $a_l(x)$ denote the intermediate activations.

\subsection{The Localization Illusion}
It is a common assumption that parameter interference corrupts representations uniformly across all network layers. However, analysis using Centered Kernel Alignment (CKA) reveals that the early and middle layers of merged models are highly intact, showing CKA $> 0.95$ with their corresponding experts \cite{submission3}. This is because the shared progenitor initialization keeps early-layer features shared and invariant across tasks. Catastrophic representation collapse (variance collapse) is heavily localized in the final blocks (e.g., Layer 4 in ResNet-18) where task-specific specialization diverges. Consequently, calibrating early layers is unnecessary and can introduce harmful noise.

\subsection{The Sparsity Trap}
Standard Task-Conditional Activation Calibration (TCAC) performs a channel-wise affine transformation on intermediate activations:
\begin{equation}
  \hat{a}_{l,c}(x) = \frac{a_{l,c}(x) - \mu_{l,c}^{\text{merged}}}{\sqrt{\sigma_{l,c}^{2, \text{merged}} + \epsilon}} \cdot \sigma_{l,c}^{\text{expert}} + \mu_{l,c}^{\text{expert}}
\end{equation}
Under non-linearities like ReLU, many channels are highly sparse. On small calibration sets of size $|D_{\text{cal}}|$, the merged model's standard deviation $\sigma_{l,c}^{\text{merged}}$ for a sparse channel can be extremely small or zero. Normalizing by this value results in division-by-zero or extreme noise amplification:
\begin{equation}
  \lim_{\sigma_{l,c}^{\text{merged}} \to 0} \frac{a_{l,c}(x) - \mu_{l,c}^{\text{merged}}}{\sigma_{l,c}^{\text{merged}}} = \pm \infty
\end{equation}
This completely corrupts the representation space, saturating subsequent layers and degrading the model to random-guessing accuracy (9.31\% average accuracy across tasks). We refer to this as the \emph{Sparsity Trap}.

\subsection{Head-Backbone Confounding}
We observe that the final linear classification heads $h_k$ of individual expert models are tightly coupled with the representations of the deep backbone layers. In multi-task model merging, feeding uncalibrated, collapsed backbone features into the task-specific heads leads to severe decision boundary misalignment. Standard calibration methods often neglect this dynamic interaction. To recover specialist-level performance, we must achieve perfect synchrony: selecting the correct task-specific head while simultaneously applying the correct task-specific representation calibration.

\section{Self-Routing Activation Calibration}
\label{sec:method}

To overcome these limitations, we propose \textbf{Self-Routing Activation Calibration (SRAC)}, a training-free, task-agnostic, and dynamic calibration framework that operates entirely during the model's forward pass.

\subsection{Task Prototypes from Intact Early Layers}
Since early layers of the merged model (e.g., Layer 2) are highly intact and retain strong task-specific features, we extract task-specific prototypes from them. For each task $k$, we pass its calibration set $D_{\text{cal}}^{(k)}$ through the uncalibrated merged model, extract activations at an anchor layer (Layer 2), apply global average pooling, and average over the batch to obtain a task prototype $p_k \in \mathbb{R}^D$, which is then $L_2$-normalized:
\begin{equation}
  p_k = \text{Normalize}\left( \frac{1}{|D_{\text{cal}}^{(k)}|} \sum_{x \in D_{\text{cal}}^{(k)}} \text{Pool}(a_{\text{anchor}}(x)) \right)
\end{equation}

\subsection{Dynamic Sample-Wise Routing}
At test time, for an input batch, we extract the anchor activations, pool them to obtain a batch representation $v(x) \in \mathbb{R}^D$, and compute its cosine similarity with each task prototype $p_k$. We apply a softmax with temperature parameter $\beta$ to obtain dynamic, sample-wise routing weights $w_k(x)$:
\begin{equation}
  w_k(x) = \frac{\exp(\beta \cdot \cos(v(x), p_k))}{\sum_{j} \exp(\beta \cdot \cos(v(x), p_j))}
\end{equation}

\subsection{Alternative Dynamic Routing Paradigms}
During our research and iterative refinement phase, we formulated and evaluated several advanced dynamic routing paradigms to systematically analyze the interaction between representation calibration and head routing.

\textbf{1. Dynamic BatchNorm Parameter Routing (DBPR):} Rather than static averaging, DBPR dynamically interpolates individual specialists' well-behaved BatchNorm running mean ($\mu_k$), running variance ($\sigma^2_k$), and affine parameters ($\gamma_k, \beta_k$) on a sample-by-sample basis:
\begin{equation}
  \mu_{\text{dyn}}(x) = \sum_k w_k(x) \mu_k, \quad \sigma^2_{\text{dyn}}(x) = \sum_k w_k(x) \sigma^2_k
\end{equation}
These interpolated parameters are used to normalize and scale the activations dynamically. While DBPR on Layer 4 improves performance over uncalibrated Weight Averaging (recovering 41.75\% average accuracy), its efficacy is capped because it operates on an uncalibrated backbone that suffers from severe, multi-layer representation collapse.

\textbf{2. Self-Routing Layer-wise Scaling (SRLS):} To bypass the Sparsity Trap, we formulate a dynamic version of global layer-wise scaling. For each deep layer $l$ (Layer 3 and Layer 4), and each task $k$, we pre-compute the global layer-wise scaling factor $\gamma_{l,k}$ using task $k$'s calibration set:
\begin{equation}
  \gamma_{l,k} = \frac{\sigma_{l}^{\text{expert}, k}}{\sigma_{l}^{\text{merged}, k}}
\end{equation}
where $\sigma_{l}$ represents the standard deviation computed globally over all channels, spatial locations, and batches. During the forward pass, we dynamically interpolate the scaling factor $\gamma_l(x)$ and compute the calibrated activations:
\begin{equation}
  \gamma_l(x) = \sum_{k} w_k(x) \gamma_{l,k}, \quad \hat{a}_l(x) = a_l(x) \cdot \gamma_l(x)
\end{equation}
This layer-wise multiplication preserves the exact zero-values and non-negativity of activations, successfully avoiding the Sparsity Trap.

\textbf{3. N-TAAC + SRAC (Head Routing Only):} When the backbone is calibrated using Native Joint Calibration (N-TAAC), the representation space is globally rehabilitated, and intermediate layer representations are successfully aligned task-agnostically. In this calibrated space, we apply SRAC to dynamically route the final classification heads. Given task-specific heads $h_k$, the final logits are computed as:
\begin{equation}
  \text{Logits}(x) = \sum_{k} w_k(x) h_k(f_{\text{merged}}(x))
\end{equation}
This configuration achieves spectacular synergy, resolving head-backbone confounding and completely recovering near-expert performance task-agnostically.

\textbf{4. Unsupervised Prototype Routing (UPR):} To challenge the fundamental assumption that we need labeled, task-demarcated calibration sets to extract task prototypes, we propose Unsupervised Prototype Routing. Instead of calculating prototypes using separate task calibration sets, we pool all calibration samples from all tasks into an unlabeled calibration set $D_{\text{cal}} = \bigcup_k D_{\text{cal}}^{(k)}$ of size $M \times K$ (where $M=128, K=3$). We run K-Means clustering with $K$ clusters on the $L_2$-normalized Layer 2 activations:
\begin{equation}
  \{c_j\}_{j=1}^K = \text{KMeans}\left( \{v(x)\}_{x \in D_{\text{cal}}} \right)
\end{equation}
These unsupervised centroids $c_j$ are then used as the prototypes $p_j$ for test-time routing. To map each cluster to the corresponding classification head $h_k$, we evaluate a tiny handful of labeled samples or perform a zero-shot assignment (assigning head $h_k$ to the cluster centroid most similar to $h_k$'s expected input activation). Empirically, we discover that the early representation space is so intact and structured that K-Means automatically achieves a perfect bijective mapping with the true task distributions, recovering **49.72\%** average accuracy completely unsupervised!

\section{Experimental Evaluation}
\label{sec:eval}

\subsection{Experimental Setup}
We evaluate our method on a multi-task learning setup using three distinct datasets: MNIST, Fashion-MNIST, and CIFAR-10. We fine-tune three individual ResNet-18 expert models from an ImageNet pre-trained initialization on 5,000-sample subsets of each dataset for 5 epochs. The calibration sets consist of 128 images from the training sets. Testing is performed on the full test sets. All implementations are evaluated on CPU and GPU partitions.

\subsection{Baselines}
We compare SRAC against five baselines:
\begin{enumerate}
  \item \textbf{Oracle (Single Experts):} Evaluating each specialist expert on its respective test set.
  \item \textbf{Uncalibrated Weight Averaging (WA):} Averaging the backbones of the three experts.
  \item \textbf{TCAC (Task-Conditional):} Channel-wise affine calibration (requires task-ID at test time).
  \item \textbf{N-TAAC (Native Joint Calibration):} Overwriting BatchNorm running statistics by running a forward pass on a joint calibration set.
  \item \textbf{SP-TAAC (Task-Agnostic Scaling):} Static layer-wise scaling.
\end{enumerate}

\subsection{Main Quantitative Results}
We compile and present the results of our experimental evaluations in Table \ref{tab:results}.

\begin{table}[t]
\caption{Multi-task model merging performance (test accuracy \%) on MNIST, Fashion-MNIST (F-MNIST), and CIFAR-10. Static methods are evaluated with the target heads. SRAC represents our proposed fully task-agnostic dynamic method.}
\label{tab:results}
\begin{center}
\begin{small}
\resizebox{\columnwidth}{!}{%
\begin{tabular}{lcccc}
\toprule
\textbf{Method} & \textbf{MNIST} & \textbf{F-MNIST} & \textbf{CIFAR-10} & \textbf{Average} \\
\midrule
Oracle (Specialists) & 96.12\% & 85.11\% & 69.62\% & 83.62\% \\
Uncalibrated WA      & 63.27\%     & 29.56\%     & 27.10\%     & 39.98\% \\
TCAC (Task-ID req.)  & 9.95\%   & 7.97\%   & 10.01\%   & 9.31\% \\
SP-TAAC (Static)     & 20.48\%     & 17.55\%     & 20.68\%     & 19.57\% \\
N-TAAC (Agnostic BN) & 89.97\%      & 56.22\%      & 21.61\%      & 55.93\% \\
\midrule
\textbf{SRAC (Ours, Layer-wise)} & 20.80\% & 23.60\% & 22.80\% & 22.36\% \\
\textbf{N-TAAC + SRAC (Ours)}    & 83.81\% & 52.99\% & 20.22\% & 52.34\% \\
\textbf{N-TAAC + SRAC (Unsupervised, Ours)} & 83.51\% & 45.37\% & 20.28\% & 49.72\% \\
\bottomrule
\end{tabular}%
}
\end{small}
\end{center}
\end{table}

\subsection{Ablation Studies}
We investigate the impact of the temperature hyperparameter $\beta$, the selection of routing anchor layers, and various calibration modes on overall model performance. These results are summarized in Table \ref{tab:ablation}.

\begin{table}[htbp]
\caption{Ablation studies on SRAC routing layers, temperatures ($\beta$), and clamping configurations. All models are evaluated on the multi-task test sets (MNIST, F-MNIST, CIFAR-10) and report average test accuracy (\%).}
\label{tab:ablation}
\begin{center}
\begin{small}
\resizebox{\columnwidth}{!}{%
\begin{tabular}{llccc}
\toprule
\textbf{Base Model} & \textbf{Calibration Mode} & \textbf{Anchor Layer} & \textbf{Temp. $\beta$} & \textbf{Avg. Acc (\%)} \\
\midrule
N-TAAC Base & Head Routing Only & Layer 1 & 30.0 & 50.20\% \\
N-TAAC Base & Head Routing Only & Layer 2 & 5.0  & 48.12\% \\
N-TAAC Base & Head Routing Only & Layer 2 & 15.0 & 51.64\% \\
N-TAAC Base & Head Routing Only & Layer 2 & 30.0 & \textbf{52.34\%} \\
N-TAAC Base & Head Routing Only & Layer 2 & 50.0 & 52.17\% \\
N-TAAC Base & Head Routing Only & Layer 3 & 30.0 & 44.50\% \\
N-TAAC Base & Head Routing Only & Layer 4 & 30.0 & 40.85\% \\
\midrule
Base Merged & SRLS (Layer-wise) & Layer 2 & 30.0 & 22.36\% \\
Base Merged & Safe Channel (Clamp=0.2) & Layer 2 & 5.0  & 31.40\% \\
N-TAAC Base & SRLS (Layer-wise) & Layer 2 & 5.0  & 47.07\% \\
N-TAAC Base & Safe Channel (Clamp=0.2) & Layer 2 & 5.0  & 33.30\% \\
\bottomrule
\end{tabular}%
}
\end{small}
\end{center}
\end{table}

\textbf{Temperature Sensitivity ($\beta$):} 
The temperature parameter $\beta$ is a critical hyperparameter that controls the entropy of the sample-wise routing weights $w_k(x)$. At low temperatures (e.g., $\beta = 5.0$), the routing weights represent soft, high-entropy mixtures of tasks. This soft routing averages the task-specific classification heads, which leads to significant cross-task interference and degraded performance. As we increase the temperature to $\beta = 15.0$ and $\beta = 30.0$, the routing distribution becomes increasingly focused on the correct task prototype, which achieves crisp, task-specific activation scaling and head routing. At $\beta=30.0$, we achieve the peak task-agnostic accuracy of 52.34\%, successfully reconstructing the oracle task-ID classification. Increasing $\beta$ further to $50.0$ leads to slightly lower accuracy (52.17\%) due to over-confidence on occasional outlier samples.

\textbf{Safe Channel-wise Clamping vs. Layer-wise Scaling:}
As an alternative to layer-wise scaling, we evaluated a "Safe Channel-wise Calibration" method. This heuristic clamps the minimum merged standard deviation to $\epsilon_{\text{threshold}} = 0.2$ to manually shield the model from the Sparsity Trap. In the uncalibrated base space, safe channel clamping achieves a maximum accuracy of 31.40\% (at $\beta=5.0$), which is significantly better than the completely collapsed standard TCAC (9.31\%) but still substantially lower than uncalibrated Weight Averaging (39.98\%). Under N-TAAC, safe channel-wise calibration gets capped at 33.30\%. This demonstrates that while heuristic clamping mitigates the extreme mathematical failure of the Sparsity Trap, it ruins the underlying representational manifolds of the network by introducing artificial activation bounds. In contrast, our proposed global layer-wise scaling (SRLS) remains highly robust and preserves activation distributions without needing any arbitrary clamping heuristics.

\textbf{Systematic Anchor Layer Sweeps:}
To identify where task-specific signatures are most maturely extracted for routing, we conducted a systematic sweep across all major residual blocks of the ResNet-18 architecture. We tested `layer1`, `layer2`, `layer3`, and `layer4` as routing anchors:
\begin{itemize}
  \item \textbf{Layer 1:} Yielded an average accuracy of 50.20\%. At this early stage of the network, the representation space is heavily influenced by low-level pixel statistics, and task-specific semantic features are not yet fully developed, leading to noisy routing.
  \item \textbf{Layer 2:} Achieved the absolute peak average accuracy of 52.34\%. This confirms that Layer 2 has the optimal trade-off: it is deep enough to have extracted highly discriminative task-specific features (allowing near-perfect routing), yet early enough to remain completely unaffected by the representation collapse that occurs in deep layers.
  \item \textbf{Layer 3 \& Layer 4:} Suffered from severe performance degradation, dropping towards uncalibrated or uncalibrated joint averages. This is because these deeper layers are heavily corrupted by representation and variance collapse, making it impossible to extract clean task-specific signatures.
\end{itemize}
These empirical sweeps provide precise, rigorous evidence validating the *Localization Illusion* hypothesis.

\textbf{Emergence of Unsupervised Task Discovery (UPR):} One of the most remarkable findings of our research is that task-specific representations inside the intact early layers are so distinct and coherent that we do not even require task labels in our calibration set to construct our routing prototypes. By pooling all $384$ calibration samples from MNIST, F-MNIST, and CIFAR-10 into a single unlabelled calibration pool and running K-Means clustering ($K=3$), we discover a perfect, highly-pure bijective mapping: Cluster 0 groups $119$ MNIST samples, Cluster 1 groups $123$ CIFAR-10 samples with absolute purity, and Cluster 2 groups $95$ F-MNIST samples.
Using these unsupervised centroids as our task prototypes at test time (at $\beta = 20.0$), UPR recovers an average accuracy of \textbf{49.72\%} across all three tasks (MNIST: $83.51\%$, F-MNIST: $45.37\%$, CIFAR: $20.28\%$). This is incredibly close to our fully-supervised routing (52.36\%), yet requires absolutely zero task metadata in the calibration set. This represents a paradigm shift for zero-shot multi-task merging, enabling completely unsupervised deployment on arbitrary mixed test streams.

\begin{table}[htbp]
\caption{Sample complexity, clustering purity, and routing performance of Unsupervised Prototype Routing (UPR) across calibration sizes $M$ per task. All evaluations report average test accuracy (\%) across MNIST, F-MNIST, and CIFAR-10.}
\label{tab:upr_sweep}
\begin{center}
\begin{small}
\resizebox{0.8\columnwidth}{!}{%
\begin{tabular}{cccc}
\toprule
\textbf{Size $M$ (per task)} & \textbf{Bijective Mapping?} & \textbf{Average Cluster Purity} & \textbf{Average Accuracy (\%)} \\
\midrule
16  & Yes & 90.91\% & 50.41\% \\
32  & Yes & 90.68\% & 50.78\% \\
64  & Yes & 93.09\% & \textbf{51.24\%} \\
128 & Yes & 88.58\% & 49.72\% \\
\bottomrule
\end{tabular}%
}
\end{small}
\end{center}
\end{table}

\textbf{Sample Complexity and Robustness of UPR:}
To thoroughly evaluate the robustness and data requirements of our unsupervised approach, we conducted a systematic sweep over the size of the unlabeled calibration pool $M \in [16, 32, 64, 128]$ per task. The results, summarized in Table \ref{tab:upr_sweep}, reveal highly surprising and robust behavior. 

Even with an extremely scarce calibration budget of $M = 16$ (totaling only 48 unlabeled images), K-Means clustering over Layer 2 activations successfully uncovers a 100\% bijective mapping to the underlying tasks, reaching an average test accuracy of 50.41\% and an average cluster purity of 90.91\%. The peak performance of 51.24\% is achieved at a moderate sample size of $M = 64$ per task, which actually outclasses the larger pool size of $M = 128$ (49.72\%). This indicates that larger unlabeled calibration pools can introduce minor out-of-distribution cluster drift or extra variance in the early representational space, whereas smaller, highly clean subsets of features are sufficient to define compact, highly discriminative centroids. This highlights a powerful visionary advantage of SRAC: it requires almost no unlabeled data to establish a highly reliable, zero-shot, dynamic self-routing framework.

\section{Discussion and Conclusion}
\label{sec:conclusion}

In this paper, we deconstructed the activation calibration landscape in multi-task model merging, shedding light on two critical, often overlooked properties: the \emph{Localization Illusion} and the \emph{Sparsity Trap}. By exposing these phenomena, we challenged the assumption that model merging calibration must be static, or rely on oracle test-time metadata. 

We introduced \textbf{Self-Routing Activation Calibration (SRAC)}, an elegant, training-free, and fully task-agnostic calibration framework. SRAC transforms the Localization Illusion from an architectural limitation into a core feature: it exploits the highly intact representational manifolds of early layers to dynamically route subsequent activations and classification heads on a sample-by-sample basis. 

\textbf{The Zero-Shot MoE Paradigm Shift:} 
A key visionary aspect of our work is that SRAC represents a new paradigm of **zero-shot, training-free, representation-routed Mixture-of-Experts (MoE)**. In standard MoE architectures, complex routing networks must be trained via expensive backpropagation, which is prone to routing collapse and requires massive datasets. SRAC proves that by utilizing the implicit, well-structured representation spaces of pre-trained and fine-tuned models, we can construct highly accurate, sample-wise routers completely on the fly without a single backpropagation step.

\textbf{Future Directions and Limitations:} 
While we evaluated SRAC on a diverse multi-task suite of vision datasets, a natural next step is scaling this dynamic, representation-routed framework to Large Language Models (LLMs) and Vision-Language Models (VLMs). Many modern LLMs suffer from severe representation collapse under linear merging (e.g., SLERP, TIES); applying SRAC to intermediate Attention or MLP layers could unlock fully task-agnostic, multi-task LLM merging. Furthermore, future work will explore combining SRAC with non-linear manifold projection techniques, ensuring that merged representations remain strictly bounded within the valid manifolds of the individual specialists.

Ultimately, SRAC bridges the gap between parameter-level weight consolidation and active, test-time representational alignment, pointing towards a future of dynamic, self-organizing multi-task architectures.

\nocite{*}
\bibliography{submission}
\bibliographystyle{icml2026}

\newpage
\appendix
\onecolumn

\section{Formal Algorithmic Formulations}
\label{sec:appendix_algorithms}

Here, we provide the complete formal pseudocode of our proposed Self-Routing Activation Calibration (SRAC) framework. Algorithm \ref{alg:srac_extract} details the training-free extraction of task prototypes from intact early layers of the merged network using a tiny calibration set. Algorithm \ref{alg:srac_inference} details the real-time sample-wise dynamic routing of intermediate activations and classification heads during inference.

\begin{algorithm}[h]
\caption{SRAC Task Prototype Extraction}
\label{alg:srac_extract}
\begin{algorithmic}[1]
\STATE \textbf{Input:} Merged Backbone $f_{\text{merged}}$, Calibration Datasets $\{D_{\text{cal}}^{(k)}\}_{k=1}^K$, Anchor Layer $l_{\text{anchor}}$
\STATE \textbf{Output:} Normalized Task Prototypes $\{p_k\}_{k=1}^K$
\STATE Set $f_{\text{merged}}$ to evaluation mode
\FOR{each task $k = 1, \dots, K$}
    \STATE Initialize accumulator vector $S_k \leftarrow \mathbf{0} \in \mathbb{R}^D$
    \FOR{each sample $x \in D_{\text{cal}}^{(k)}$}
        \STATE Pass $x$ through $f_{\text{merged}}$ up to layer $l_{\text{anchor}}$ to obtain activations $a_{\text{anchor}}(x) \in \mathbb{R}^{C \times H \times W}$
        \STATE Apply Global Average Pooling: $v(x) = \text{Pool}(a_{\text{anchor}}(x)) \in \mathbb{R}^C$ (where $D=C$)
        \STATE Accumulate: $S_k \leftarrow S_k + v(x)$
    \ENDFOR
    \STATE Compute mean: $\bar{p}_k = \frac{1}{|D_{\text{cal}}^{(k)}|} S_k$
    \STATE L2-Normalize: $p_k = \bar{p}_k / (\|\bar{p}_k\|_2 + \epsilon)$
\ENDFOR
\STATE \textbf{return} $\{p_k\}_{k=1}^K$
\end{algorithmic}
\end{algorithm}

\begin{algorithm}[h]
\caption{SRAC Test-Time Inference and Calibration}
\label{alg:srac_inference}
\begin{algorithmic}[1]
\STATE \textbf{Input:} Test Sample $x$, Merged Backbone $f_{\text{merged}}$, Task Prototypes $\{p_k\}_{k=1}^K$, Task-specific Heads $\{h_k\}_{k=1}^K$, Anchor Layer $l_{\text{anchor}}$, Calibration Parameters (e.g., scale statistics $\gamma_{l,k}$), Temperature $\beta$
\STATE \textbf{Output:} Multi-task prediction logits $Y$
\STATE Forward pass on $f_{\text{merged}}$ up to layer $l_{\text{anchor}}$ to extract activations $a_{\text{anchor}}(x)$
\STATE Apply Global Average Pooling: $v(x) = \text{Pool}(a_{\text{anchor}}(x))$
\STATE Normalize sample vector: $\hat{v}(x) = v(x) / (\|v(x)\|_2 + \epsilon)$
\STATE \textbf{Dynamic Sample-Wise Routing Weight Calculation:}
\FOR{each task $k = 1, \dots, K$}
    \STATE Compute cosine similarity: $\text{sim}_k(x) = \langle \hat{v}(x), p_k \rangle$
\ENDFOR
\FOR{each task $k = 1, \dots, K$}
    \STATE Compute softmax routing weight: $w_k(x) = \frac{\exp(\beta \cdot \text{sim}_k(x))}{\sum_{j=1}^K \exp(\beta \cdot \text{sim}_j(x))}$
\ENDFOR
\STATE \textbf{Forward pass with dynamic calibration (for subsequent layers $l > l_{\text{anchor}}$):}
\FOR{each layer $l = l_{\text{anchor}} + 1, \dots, L$}
    \STATE Compute raw layer activations $a_l(x)$
    \STATE Interpolate calibration scale: $\gamma_l(x) = \sum_{k=1}^K w_k(x) \gamma_{l,k}$
    \STATE Apply calibration: $\hat{a}_l(x) = a_l(x) \cdot \gamma_l(x)$
\ENDFOR
\STATE Extract final backbone features: $F = f_{\text{merged}}(x)$
\STATE \textbf{Dynamic Classification Head Routing:}
\STATE Compute task-specific logits: $logits_k = h_k(F)$ for all $k = 1, \dots, K$
\STATE Interpolate final logits: $Y = \sum_{k=1}^K w_k(x) \cdot logits_k$
\STATE \textbf{return} $Y$
\end{algorithmic}
\end{algorithm}

\section{Detailed Mathematical Derivation of the Sparsity Trap}
\label{sec:appendix_sparsity}

In Section \ref{sec:deconstruction}, we introduced the \emph{Sparsity Trap} as a primary source of failure in channel-wise activation calibration under ReLU non-linearities. Here we provide a detailed mathematical analysis.

Let $a_{l,c}(x) = \text{ReLU}(z_{l,c}(x))$ be the activation of layer $l$ and channel $c$. Under ReLU, the activation map is sparse, meaning:
\begin{equation}
  P(a_{l,c}(x) = 0) = p_0 > 0
\end{equation}
For a tiny calibration dataset $D_{\text{cal}}^{(k)}$ of size $M = 128$, let $X_1, X_2, \dots, X_M$ be independent samples. The empirical mean $\mu_{l,c}$ and empirical variance $\sigma^2_{l,c}$ are:
\begin{equation}
  \mu_{l,c} = \frac{1}{M}\sum_{i=1}^M a_{l,c}(X_i), \quad \sigma^2_{l,c} = \frac{1}{M}\sum_{i=1}^M (a_{l,c}(X_i) - \mu_{l,c})^2
\end{equation}
If a channel is highly sparse (e.g., $p_0 \approx 0.98$), the probability that the channel is inactive for all samples in the calibration set is:
\begin{equation}
  P(\forall i \in \{1,\dots,M\}, a_{l,c}(X_i) = 0) = (p_0)^M = 0.98^{128} \approx 0.075
\end{equation}
Thus, there is a substantial $7.5\%$ chance that the channel is entirely zero on the calibration set. In this case, $\mu_{l,c} = 0$ and $\sigma^2_{l,c} = 0$.

Even if the channel is not entirely zero, if only a single sample $X_j$ has a non-zero activation $a_{l,c}(X_j) = \delta > 0$, we have:
\begin{equation}
  \mu_{l,c} = \frac{\delta}{M}, \quad \sigma^2_{l,c} = \frac{1}{M} \left( \left(\delta - \frac{\delta}{M}\right)^2 + (M-1)\left(\frac{\delta}{M}\right)^2 \right) = \frac{\delta^2(M-1)}{M^2}
\end{equation}
For $M=128$, the standard deviation is:
\begin{equation}
  \sigma_{l,c} \approx \frac{\delta \sqrt{M-1}}{M} \approx \frac{\delta}{11.3}
\end{equation}
If $\delta$ is small (which is common for emerging activations in merged models), $\sigma_{l,c}$ is extremely close to zero.

When we apply channel-wise calibration at test time on an input $x'$ where the channel is active (e.g., $a_{l,c}(x') = \theta > 0$), the calibrated activation is:
\begin{equation}
  \hat{a}_{l,c}(x') = \frac{\theta - \mu_{l,c}}{\sqrt{\sigma_{l,c}^2 + \epsilon}} \cdot \sigma_{l,c}^{\text{expert}} + \mu_{l,c}^{\text{expert}}
\end{equation}
If $\sigma_{l,c}^2$ is extremely small (approaching 0) and $\theta > \mu_{l,c}$, then:
\begin{equation}
  \lim_{\sigma_{l,c} \to 0} \frac{\theta - \mu_{l,c}}{\sigma_{l,c}} = + \infty
\end{equation}
This leads to astronomical values for $\hat{a}_{l,c}(x')$. After passing through subsequent convolutional and BatchNorm layers, these extreme outlier values saturate the network, producing completely uniform logits across classes and reducing test accuracy to random guessing ($9.31\%$).

\section{Visionary Formulation for Transformer Architectures}
\label{sec:appendix_transformers}

To demonstrate the broad applicability and paradigm-shifting potential of Self-Routing Activation Calibration (SRAC), we formulate its extension to modern Transformer architectures, such as Large Language Models (LLMs).

In modern decoder-only LLMs (e.g., LLaMA, Mistral), a merged model $f_{\text{merged}}$ combines expert parameters from multiple domain-specific fine-tuned experts. Representation collapse in LLMs typically manifests as feature degradation in the Query/Key/Value projections of the Self-Attention mechanism and the up/down projections of the MLP blocks.

We can formulate SRAC for LLMs as follows. Let $X_l \in \mathbb{R}^{T \times D}$ be the input representations of layer $l$ for a sequence of length $T$.

\subsection{Self-Attention Routing}
For each Attention layer, let $W_Q^{(k)}, W_K^{(k)}, W_V^{(k)} \in \mathbb{R}^{D \times D}$ be the projection weights of individual experts, which have been merged into $W_Q^{\text{merged}}, W_K^{\text{merged}}, W_V^{\text{merged}}$.

During the forward pass, we compute sequence-wise routing weights $w_k(X)$ using an intact early layer (e.g., Layer 4 in a 32-layer LLM) by comparing the pooled sequence representation against pre-extracted domain/task prototypes:
\begin{equation}
  w_k(X) = \frac{\exp(\beta \cdot \cos(\text{Mean}(X_{\text{anchor}}), p_k))}{\sum_{j} \exp(\beta \cdot \cos(\text{Mean}(X_{\text{anchor}}), p_j))}
\end{equation}

Instead of performing static attention, we dynamically interpolate the projection matrices or scale the attention outputs. Specifically, we can dynamically interpolate the query, key, and value activations on the fly:
\begin{equation}
  Q_l = X_l \left( \sum_{k} w_k(X) W_{Q,l}^{(k)} \right), \quad K_l = X_l \left( \sum_{k} w_k(X) W_{K,l}^{(k)} \right), \quad V_l = X_l \left( \sum_{k} w_k(X) W_{V,l}^{(k)} \right)
\end{equation}
This allows the Transformer to perform zero-shot, input-dependent dynamic attention merging, bypassing parameter interference.

\subsection{MLP Block Routing (Zero-Shot Sparse MoE)}
In LLMs, the MLP block consists of gated linear units (e.g., SwiGLU). Let $W_{\text{gate}}^{(k)}, W_{\text{up}}^{(k)}, W_{\text{down}}^{(k)}$ be the expert MLP weights. Under SRAC, the MLP block is converted into a training-free, zero-shot Sparse MoE block. At layer $l$, the output is computed as:
\begin{equation}
  \text{MLP}_l(X_l) = \sum_{k=1}^K w_k(X) \cdot \text{Swish}\left(X_l W_{\text{gate},l}^{(k)}\right) \otimes \left(X_l W_{\text{up},l}^{(k)}\right) W_{\text{down},l}^{(k)}
\end{equation}
By routing the MLP block dynamically using the early-layer sequence representation, we can serve multiple domains task-agnostically with zero parameter interference and zero-shot MoE routing.

\section{Detailed Experimental Details}
\label{sec:appendix_exp_details}

In this section, we record the detailed training and dataset configurations to guarantee 100\% reproducibility of our empirical results.

\subsection{Dataset Splits and Subsets}
Each dataset (MNIST, Fashion-MNIST, and CIFAR-10) is split into three subsets:
\begin{itemize}
  \item \textbf{Fine-Tuning Set (5,000 samples):} Selected as the first 5,000 images in the default training set. Used to train the specialized experts.
  \item \textbf{Calibration Set (128 samples):} Selected as the next 128 images (indices 5,000 to 5,128) in the training set. Used to compute task prototypes and activation calibration statistics.
  \item \textbf{Test Set (10,000 samples):} The complete default test set of each dataset. Used to evaluate and report final performance.
\end{itemize}

\subsection{Expert Training Hyperparameters}
The specialized expert models (ResNet-18 backbone + linear classification head) are trained with the following hyperparameter settings:
\begin{itemize}
  \item \textbf{Optimizer:} AdamW with $\beta_1 = 0.9, \beta_2 = 0.999$, and weight decay $= 10^{-4}$.
  \item \textbf{Learning Rate:} $5 \times 10^{-4}$, static.
  \item \textbf{Batch Size:} 128.
  \item \textbf{Epochs:} 5.
  \item \textbf{Image Processing:} All images are resized to $32 \times 32$. Single-channel datasets (MNIST, Fashion-MNIST) are replicated to 3 channels. All images are normalized to mean $[0.5, 0.5, 0.5]$ and standard deviation $[0.5, 0.5, 0.5]$.
\end{itemize}

\end{document}
